\section{RNN 的应用场景}

RNN 不仅可以用作语言模型，还广泛应用于各种序列处理任务。

\subsection{序列标注}

RNN 可以在每个时间步产生一个输出，用于\textbf{序列标注}（Sequence Tagging）任务：

\begin{figure}[H]
\centering
\includegraphics[width=0.95\textwidth]{slides-images/slide-065.jpg}
\caption{RNN 用于词性标注（POS Tagging）：输入一个句子，为每个词预测其词性标签（DT, JJ, NN, VBN 等）\protect\footnotemark}
\end{figure}
\footnotetext{Slides 第66页。}

典型应用包括：
\begin{itemize}
    \item \textbf{词性标注}（Part-of-Speech Tagging）：为每个词标注名词、动词、形容词等语法类别
    \item \textbf{命名实体识别}（Named Entity Recognition, NER）：识别人名、地名、组织名等
\end{itemize}

\subsection{文本分类与情感分析}

对于需要对整个序列产生单一判断的任务，可以使用 RNN 最后一个时间步的隐藏状态（或所有隐藏状态的某种聚合）作为序列的表示，再接分类层。

\begin{figure}[H]
\centering
\includegraphics[width=0.95\textwidth]{slides-images/slide-066.jpg}
\caption{RNN 用于情感分类：处理完整个句子后，用最终隐藏状态进行分类\protect\footnotemark}
\end{figure}
\footnotetext{Slides 第67页。}

\subsection{条件语言模型与编码器-解码器架构}

通过用不同的信息初始化 RNN 的隐藏状态，可以构建\textbf{条件语言模型}：

\begin{figure}[H]
\centering
\includegraphics[width=0.95\textwidth]{slides-images/slide-070.jpg}
\caption{条件语言模型：以音频特征初始化隐藏状态，RNN 生成对应的文字转录。这是语音识别和机器翻译的基础框架\protect\footnotemark}
\end{figure}
\footnotetext{Slides 第71页。}

\begin{knowledgebox}{从 RNN 到编码器-解码器}
条件语言模型的思想自然延伸为\textbf{编码器-解码器}（Encoder-Decoder）架构：一个 RNN（编码器）读取输入序列并产生一个固定长度的隐藏表示，另一个 RNN（解码器）以该表示为条件生成输出序列。这一架构在机器翻译（Seq2Seq）中取得了巨大成功，将在后续课程中详细讨论。
\end{knowledgebox}

\subsection{本章小结}

\begin{itemize}
    \item RNN 是一种通用的序列处理架构，适用于多种 NLP 任务
    \item 每步都有输出 $\rightarrow$ 序列标注（词性标注、命名实体识别）
    \item 只看最终输出 $\rightarrow$ 文本分类（情感分析）
    \item 用外部信息初始化 $\rightarrow$ 条件生成（语音识别、机器翻译）
\end{itemize}

%% ========== 总结与延伸 ==========

